% !TeX program = xelatex
\documentclass[11pt]{impeart}

% The manual deliberately applies a global house style.  Ordinary documents
% should normally use the registered automatic behavior of \UseFont or the template `fonts` key.
\UseFont{libertinus}[global]
\setmonofont[
  BoldFont=lmmonolt10-bold.otf,
  ItalicFont=lmmono10-italic.otf,
  BoldItalicFont=lmmonolt10-boldoblique.otf
]{lmmono10-regular.otf}
\UseFeatures{hyperlinks,tables}

\usepackage{catchfile}
\usepackage{metalogo}
\usepackage{pdfpages}

\newcommand{\IMPEVersionFile}{../../VERSION}
\newcommand{\IMPEShowcaseFile}{../showcase/impe-showcase.pdf}

\IfFileExists{\IMPEVersionFile}
  {\CatchFileEdef{\IMPEVersion}{\IMPEVersionFile}{\endlinechar=-1\relax}}
  {\PackageError{impe-manual}
     {Required repository resource ../../VERSION is missing}
     {Compile from manual/en or run scripts/build_manual.ps1 from the repository root.}}

% Keep the generated manual reproducible when SOURCE_DATE_EPOCH is fixed.
\ifdefined\XeTeXversion
  \special{pdf:trailerid [
    <5677bc84f2edee3de328d781062f390b>
    <778a1a55595b871c7ecbd03469f75deb>
  ]}
\else\ifdefined\pdftrailerid
  \expandafter\pdftrailerid\expandafter{IMPE-\IMPEVersion-manual-en}
\fi
\fi

\title{IMPE \LaTeX{} System}
\subtitle{User Manual}
\author{YANG Sikai}
\date{Version \IMPEVersion}

\newcommand{\IMPERepo}{https://github.com/KelvinYangBK67/IMPE-LaTeX-System}
\newcommand{\IMPEIssues}{https://github.com/KelvinYangBK67/IMPE-LaTeX-System/issues}
\newcommand{\IMPEShowcaseURL}{https://github.com/KelvinYangBK67/IMPE-LaTeX-System/blob/main/manual/showcase/impe-showcase.pdf}
\newcommand{\IMPETaggedShowcaseURL}{https://github.com/KelvinYangBK67/IMPE-LaTeX-System/blob/v\IMPEVersion/manual/showcase/impe-showcase.pdf}

\begin{document}

\maketitle

\begin{center}
\small English reference edition
\end{center}

\tableofcontents

\section{Introduction}

The IMPE \LaTeX{} System is a modular \XeLaTeX{} framework for reusable layouts,
font routing, multilingual typesetting, and optional document features.

IMPE addresses this by separating the system into reusable parts.  From the
user's point of view, the most important distinction is between \emph{layouts},
\emph{fonts}, and \emph{features}.  A layout describes the principal visual
organization of a document.  Font families provide global or local font
behavior, including script-aware routing where necessary.  Features provide
optional functionality such as mathematics, hyperlinks, citations, tables, or
images.  These parts can be selected independently or combined through one
unified template interface.

The normal use of IMPE is intentionally short.  A wrapper class can establish a
sensible document layout and a default global font setup, after which a user
loads the additional families or features required by the document. When
more explicit control is wanted, the same components can be declared in a
single \verb|\UseTemplateSet| call.

\XeLaTeX{} is the primary document engine for IMPE.  This choice allows the system
to rely on Unicode text, OpenType fonts, \texttt{fontspec}, and \XeTeX{}'s
inter-character mechanisms while retaining the familiar \LaTeX{} document model.
Some specialized internal routes may invoke other engines where a particular
script or renderer requires them, but ordinary IMPE documents are intended to
be compiled with \XeLaTeX{}.

This manual is a user-oriented introduction. Detailed subsystem documentation is maintained in the
repository under \texttt{docs/}.  In particular, the font, layout, feature, and
system documents describe registration fields and implementation details that
would be unnecessarily repetitive here.  This English manual is the reference
edition: translated manuals should follow its structure and technical content.

A visual survey of the system is available in the project showcase at
\href{\IMPEShowcaseURL}{the current IMPE showcase PDF}.  The edition frozen
with this manual is available from the
\href{\IMPETaggedShowcaseURL}{version tag}.  That canonical repository PDF is
also reproduced in Appendix~\ref{app:showcase}.  It is intended to demonstrate
the range of scripts, writing directions, layouts, and feature combinations
that motivated the system.

\section{Installation}

\subsection{Distribution forms}

IMPE 1.1.0 is prepared in three distribution forms. They share the same
public class/package names and the core \XeLaTeX{} runtime, but differ in
font assets and installation workflow.

\begin{description}
  \item[\texttt{Full}] The complete managed release, including the runtime,
    compatibility entry points, vendored XSR, and only those local font files
    approved for redistribution and actually present when the archive is built.
    A public Full archive is not a promise that every catalogued font is bundled.
  \item[\texttt{Core}] The managed runtime, compatibility entry points, vendored
    XSR and installers, without a local font library. Existing \TeX{} Live and
    system fonts can be resolved directly; other required files must be supplied.
  \item[\texttt{CTAN}] The \texttt{impe-framework.zip} source and runtime
    distribution rooted at \texttt{impe-framework/}. It includes canonical
    \texttt{impe*} entry points, the vendored XSR sources, and documentation
    resources, but no font binaries. CTAN installation itself does not run
    \texttt{pip} or create a Python environment.
\end{description}

The project is designed to keep the package runtime separate from licensing and
availability decisions about large third-party font collections. The final
public Full archive uses an explicit reviewed font allowlist.

\subsection{Installing Full or Core}

Direct installation places the managed package under
\texttt{TEXMFHOME/tex/latex/impe/}. Its canonical entry points are
\texttt{impe.sty}, \texttt{impeart.cls}, \texttt{impebook.cls},
\texttt{impereport.cls} and \texttt{impebeamer.cls}; the
\texttt{\_zh} wrappers and historical \texttt{next*} names are retained for
compatibility in direct releases.

The provided installers work from an extracted Full/Core release or from a
complete source checkout. Typical entry points are:

\begin{verbatim}
# Windows PowerShell, in an extracted release:
.\install.bat
# Windows PowerShell, in a source checkout:
.\scripts\install.ps1

# Unix, in an extracted release:
sh ./install.sh
# Unix, in a source checkout:
sh ./scripts/install.sh
\end{verbatim}

The Windows installer accepts \texttt{-TexmfRoot} and
\texttt{-SkipXsrPython}; the Unix installer accepts an optional TEXMF
destination as its positional argument. When its installation completes, refresh
the \TeX{} filename database if the installer could not run
\texttt{mktexlsr}. A package installed through the \TeX{} distribution's own
package manager follows that distribution's installation rules instead.

\subsection{XSR and the Python renderer}

IMPE bundles XSR \textbf{1.3}. Full and Core installers copy its \TeX{}
integration and normally set up a \emph{private} Python virtual environment for
its renderer, without installing XSR into the user's global Python environment.
Dynamic Egyptian hieroglyphic composition, native Khitan stacking, and XSR
vector glyph processing may require that renderer.

The normal managed installation requires Python \textbf{3.11 or newer},
\texttt{venv}, and \texttt{pip}; the initial dependency installation may require
network access. A \TeX{}-only installation can skip automatic Python setup:

\begin{verbatim}
# Windows:
.\scripts\install.ps1 -SkipXsrPython
# Unix:
IMPE_SKIP_XSR_PYTHON=1 sh ./scripts/install.sh
\end{verbatim}

The \TeX{}-only option installs the \TeX{} runtime, not a working dynamic
Python renderer; configure the renderer separately if using XSR. Dynamic XSR
jobs need an allowed XeLaTeX shell-escape invocation or the XSR preprocessing
route. Ordinary documents not requesting those renderers do not start Python.
The XSR script backends also need the real readable TTF/OTF files used to
produce their glyph outlines: an installed font family name by itself is
insufficient for the Python font-file reader.

\subsection{Font resources and local configuration}

For a registered face IMPE first checks the corresponding local/bundled file,
then attempts resolution through the \TeX{} font system or operating system.
Core and CTAN users do not have to copy an already resolvable system font into
the IMPE repository. Other registered fonts need the necessary files supplied
by the user, subject to their respective licenses.

The normal local font root is \texttt{assets/fonts/}. To point at a separately
managed font library, put the following in \texttt{impe.local.tex}, or issue
the command from the document preamble:

\begin{verbatim}
\SetCatalogFontRoot{D:/Fonts/IMPE}
\end{verbatim}

The path is only an example. Full installations can provide an automatically
managed local root. Unresolved ordinary fonts default to a soft warning and
LaTeX fallback; that fallback does not substitute for the actual outline file
required by an XSR script renderer.

\subsection{A quick installation check}

On a system with Libertinus available, the following checks that the wrapper,
registry and local font command can be found:

\begin{verbatim}
\documentclass{impeart}
\UseFont{libertinus}
\begin{document}
IMPE is ready. \LIB{Local Libertinus text.}
\end{document}
\end{verbatim}

This checks normal \XeLaTeX{} operation; it does not test XSR/Python or every
optional feature. See the project README and the installation checklist for
separate XSR and release validation.

\section{Quick Start}

The simplest IMPE workflow is to choose a wrapper class, select any desired
font family or feature, and write the document normally. IMPE provides a
reusable setup layer around ordinary \LaTeX{} document structure.

A small article can therefore look like this:

\begin{verbatim}
\documentclass[11pt]{impeart}

\UseFont{libertinus}
\UseFeatures{math,hyperlinks}

\title{A Sample Document}
\subtitle{Typeset with IMPE}
\author{Author Name}

\begin{document}

\maketitle

\section{Introduction}
\LIB{Hello, world.}

\end{document}
\end{verbatim}

The class already supplies the normal English article layout and base font.
The mode-free \verb|\UseFont{libertinus}| call follows that family's registered
automatic behavior, applying its global stack and making its local \verb|\LIB|
command available, while \verb|\UseFeatures| activates two optional feature
modules.  The rest is ordinary \LaTeX{}.

For a document whose setup is better expressed as a single declaration, the
same idea can be written through \verb|\UseTemplateSet|:

\begin{verbatim}
\documentclass{impeart}

\UseTemplateSet{
  fonts = {libertinus},
  features = {math,hyperlinks}
}
\end{verbatim}

A template set may also choose a layout and additional font families.  In
practice, wrapper classes already supply a suitable default layout, so most
documents state the settings that differ from those defaults.

IMPE can also be loaded as a package underneath a standard \LaTeX{} class:

\begin{verbatim}
\documentclass{article}
\usepackage{impe}

\UseTemplateSet{
  layout = en_doc,
  fonts = {cmu,libertinus},
  features = {hyperlinks}
}
\end{verbatim}

This form is useful when the underlying class must remain explicit or when every
IMPE component should be selected manually.  The wrapper classes and the raw
package use the same underlying system.


\section{Ready-to-use Templates}

The repository includes a small collection of starter projects under
\texttt{templates/}. These are intended for actual writing. Regression
fixtures and internal demonstrations live under \texttt{tests/}.

The current starter templates are:

\begin{description}
  \item[\texttt{article\_en/}] English article starter using \texttt{impeart}.
  \item[\texttt{article\_zh/}] Chinese article starter using \texttt{impeart\_zh}.
  \item[\texttt{book\_en/}] English book starter using \texttt{impebook}.
  \item[\texttt{book\_zh/}] Chinese book starter using \texttt{impebook\_zh}.
  \item[\texttt{beamer\_en/}] English presentation starter using \texttt{impebeamer}.
  \item[\texttt{beamer\_zh/}] Chinese presentation starter using \texttt{impebeamer\_zh}.
\end{description}

Each template follows the installed-package style and contains realistic
metadata, sectioning, and body structure, together with representative table,
figure, or slide material where appropriate.  A new project can therefore begin
by copying the nearest template directory and editing the document contents
using the supplied preamble as a starting point.

The templates deliberately use the canonical \texttt{impe*} entry points and
show normal user-facing practice. Development fixtures live under
\texttt{tests/}.

\section{Document Classes and Public Entry Points}

\subsection{Wrapper classes}

IMPE provides wrapper classes for the four principal document forms used by the
system:

\begin{description}
  \item[\texttt{impeart}] English-oriented article wrapper.  Its normal layout
    is \texttt{en\_doc} and its default global family is \texttt{cmu}.
  \item[\texttt{impebook}] English-oriented book wrapper.  Its normal layout
    is \texttt{en\_book} and its default global family is \texttt{cmu}.
  \item[\texttt{impereport}] English-oriented report wrapper.  It uses the
    \texttt{en\_doc} layout and \texttt{cmu} by default.
  \item[\texttt{impebeamer}] Presentation wrapper.  It selects the
    \texttt{beamer} layout and \texttt{cmu} by default.
\end{description}

Chinese-oriented counterparts are provided as \texttt{impeart\_zh},
\texttt{impebook\_zh}, \texttt{impereport\_zh}, and
\texttt{impebeamer\_zh}.  These select the corresponding Chinese document
settings and include \texttt{shanggu} in the default global font configuration
alongside \texttt{cmu}. This default configuration supports multilingual
content in either wrapper family.

Unknown class options are passed through to the underlying standard class.  A
normal option such as

\begin{verbatim}
\documentclass[12pt,twoside]{impeart}
\end{verbatim}

therefore continues to control the underlying \texttt{article} class in the
usual way.

\subsection{Using \texttt{impe.sty} directly}

The wrapper classes share one implementation and provide convenient defaults.
A user may also keep a standard class and load

\begin{verbatim}
\usepackage{impe}
\end{verbatim}

followed by an explicit layout, font, and feature configuration.  This is the
recommended route when the document class is supplied by a journal, publisher,
or another package and must remain the document's class.

\subsection{Titles and subtitles}

IMPE supports \verb|\subtitle{...}| alongside the standard \verb|\title| and
\verb|\author| commands.  Wrapper classes apply their own title-block defaults,
while the ordinary \LaTeX{} title mechanism remains recognizable.  A typical
preamble is simply

\begin{verbatim}
\title{Main Title}
\subtitle{A More Specific Description}
\author{Author Name}
\end{verbatim}

followed by \verb|\maketitle| in the document body.

\section{The Template Model}

\subsection{Four setup keys}

The principal unified command is \verb|\UseTemplateSet|.  Its public setup
model is intentionally small:

\begin{verbatim}
\UseTemplateSet{
  layout = <preset>,
  fonts = {<families>},
  globalfonts = {<families>},
  features = {<features>}
}
\end{verbatim}

The four keys have different roles.

\begin{description}
  \item[\texttt{layout}] selects a public layout preset appropriate to the
    active document class.
  \item[\texttt{fonts}] loads families through their registered automatic
    behavior.  This is the normal choice: a family may define a local command,
    apply a document-wide font, or install range-specific routing.
  \item[\texttt{globalfonts}] explicitly forces the listed families through
    their global mode or global routing path.
  \item[\texttt{features}] loads optional functional modules.
\end{description}

\texttt{mainfonts} is accepted as an alias of \texttt{globalfonts}; both keys
provide explicit global overrides; \texttt{fonts} provides the automatic interface.

The template interface provides one predictable place to express a complete
configuration. Wrapper classes already apply their own layout and global-font
defaults, so a document may specify an additional feature list alone.

\subsection{Single-purpose shortcuts}

For focused configuration, IMPE provides
public shortcuts corresponding to the individual subsystems:

\begin{verbatim}
\UseLayout{en_doc}
\UseLayouts{...}

\UseFont{libertinus}
\UseFonts{arabic,tibetan}

\UseFont{libertinus}[global]
\UseLocalFonts{arabic,tibetan} % explicit local override
\UseGlobalFonts{libertinus}    % explicit global override

\UseFeature{hyperlinks}
\UseFeatures{math,hyperlinks,tables}
\end{verbatim}

The mode-free font forms are the recommended shortcuts for ordinary use.  The
local and global aliases remain useful when a document intentionally overrides
a family's registered behavior.  The unified and single-purpose interfaces are
two views of the same public model.

\subsection{Compositional templates}

An IMPE template combines reusable components. A layout can
be reused with different font choices, the same feature set can be used in an
article or book, and a script-specific font family can be added while
preserving unrelated document functionality. This composition is the central
design principle behind the public interface.

For example, the following declarations compose two different documents
from the same template interface:

\begin{verbatim}
% English article
\UseTemplateSet{
  layout = en_doc,
  fonts = {libertinus},
  features = {math,hyperlinks}
}

% Chinese book with additional script support
\UseTemplateSet{
  layout = zh_book,
  fonts = {cmu,shanggu,sanskrit,tibetan},
  features = {citations,index,hyperlinks}
}
\end{verbatim}

The second example states the families and features it needs; the internal
font and layout subsystems remain responsible for the implementation details.

\section{Fonts, Script Routing, and Multilingual Typesetting}

IMPE provides a short preset interface for everyday documents and a separate
Writing Profile interface for explicit multilingual configuration. Both are
supported public APIs in 1.1.0; the newer model does not invalidate historical
font IDs or their local commands.

\subsection{Simple registered presets: \texttt{\string\UseFont}}

The easiest way to load an existing registered preset is
\verb|\UseFont{id}|; the \texttt{fonts} key of
\verb|\UseTemplateSet| accepts a list of the same IDs.
Omitting the mode follows that preset's \emph{automatic dispatch}, which can
enable a local command, a document-wide font binding, and/or Unicode-range
routing. It is not equivalent to simply reading the catalog's
\texttt{defaultmode} field.

\begin{verbatim}
\UseFont{libertinus}       % automatic: local command plus global candidate
\UseFonts{sanskrit,tibetan,arabic}
\UseFont{libertinus}[local]   % local declaration only
\UseFont{libertinus}[global]  % explicit global declaration
\end{verbatim}

The first call makes \verb|\LIB{...}| available and also registers
Libertinus as an explicit global Latin candidate. In a wrapper document it
can replace the implicit CMU primary according to the precedence rules below.
An explicit \texttt{[local]} confines the request to the local declaration.
An explicit \texttt{[global]} requires a global-capable preset. The
\verb|\UseLocalFont(s)|, \verb|\UseGlobalFont(s)| and
\verb|\UseMainFont(s)| aliases give the corresponding deliberate overrides.

A local command, such as \verb|\SA{...}|, applies to its text argument,
preserving the script's shaping and special renderer. The generic
\verb|\Font{preset}{text}| offers the same forced-local selection without
needing an earlier \verb|\UseFont| call. Family commands and the preset
interface remain especially convenient for established editorial styles.

\subsection{Writing Profiles and Font Resources}

The explicit writing model separates four kinds of information:
\emph{Unicode block data}, \emph{Writing Profiles} (codepoint coverage and
script-dependent behavior), \emph{font resources/packs} (real font faces), and
\emph{backward-compatible registered presets}. A profile is not a language
name, a font family or a legacy preset. For example, \texttt{sanskrit} is a
preset with historically established typography; \texttt{devanagari} is a
Writing Profile to which a suitable font resource can be attached.

Version 1.1.0 provides 69 curated profiles, 85 resource aliases and all 59
original registered preset IDs. The profile determines Unicode coverage and
the appropriate script/renderer policy; a resource override changes the font
choice, not the profile's identity. Thus \texttt{chinese},
\texttt{japanese} and \texttt{korean} profiles can share Han codepoints:
Unicode alone cannot guess which regional Han form is intended.

A global writing declaration uses:

\begin{verbatim}
\UseWriting{[shanggu]chinese}
\UseWriting{[tiro]devanagari}
\UseWriting{[naskh]arabic}
\end{verbatim}

The general syntax is \verb|\UseWriting{[pack,options]profile}|. Multiple
entries may be comma-separated outside brackets and braces. A known pack name
may appear first in brackets; \texttt{font=} selects a resolvable font family
name, whereas \texttt{file=} selects a physical file path. These mutually
exclusive resource selectors are distinct from the Writing Profile. Additional
keys include \texttt{stylemap}, face overrides (for example
\texttt{italic=} or \texttt{sansbold=}) and \texttt{writing=}.
Use a global declaration only when its script coverage and rendering mode
are appropriate. Global vertical writing is not supported.

\subsection{Named and anonymous local writing}

Use \verb|\DeclareWriting{name}{configuration}| to save a reusable
configuration without activating it globally. \verb|\Font{name}[overrides]{text}|
forces that configuration locally; \verb|\Writing{configuration}{text}| is
its anonymous counterpart. For example:

\begin{verbatim}
\DeclareWriting{han-note}
  {[shanggu,stylemap=full,italic=zhuque]chinese}
\begin{document}
\Font{han-note}{漢字}
\Font{sanskrit}{संस्कृतम्}
\Writing{[tiro]devanagari}{संस्कृतम्}
\end{document}
\end{verbatim}

Local selection applies one font resource and renderer to the
\emph{entire supplied argument}; it does not automatically subdivide that
argument according to the profile's Unicode ranges. Nested explicit local
calls can override the outer selection. A local configuration may also
specify only \texttt{font=} or \texttt{file=} when no script-specific
Profile behavior is needed.

The \texttt{stylemap=native} policy favors the resource's genuine style
faces; \texttt{stylemap=full} additionally permits the documented
cross-family stylistic mappings. Several multi-face Western packs default
to the full mapping. Style selection and missing-character recovery are
separate operations: an available Bold face may still lack a particular
character. See \texttt{docs/USEWRITING-1.0.5.md} for complete options
and profile/resource names; the filename records the API's introduction,
not its support ceiling.

\subsection{Global precedence and automatic fallback}

A wrapper's implicit CMU (and, in Chinese wrappers, Shanggu) setup supplies
the default font pool, but \emph{explicit user declarations take priority}
over those class defaults. Within one priority tier, the first declared
Latin global primary normally remains primary. A later global registration
does not automatically replace the earlier explicit primary. For CJK
cross-profile conflicts within a tier, the existing Chinese--Japanese--Korean
priority is retained. Explicitly selecting a previously implicit default can
promote it to the user tier.

For example, after \verb|\documentclass{impeart}|, calling
\verb|\UseFont{libertinus}| explicitly lets Libertinus supersede the
class-provided CMU primary. But in

\begin{verbatim}
\UseGlobalFonts{cmu,libertinus}
\end{verbatim}

the first explicit Latin primary, CMU, remains primary. This difference is
intentional. It also explains why an \emph{Auto} table entry describes
registration capabilities and effects, not an unconditional promise that
every requested font becomes the document's main face.

\emph{Automatic missing-glyph fallback currently applies to CJK only.}
The xeCJK backend searches the active resource's declared faces using the
Serif/Sans $\times$ Regular/Bold $\times$ Upright/Italic
(2$\times$2$\times$2) style matrix. It then considers available registered
resources of the same Writing Profile, followed by other CJK profiles.
The global fallback resource pool places explicit user resources before
implicit class defaults and maintains first-registration order within a tier.
Candidates with a valid face but a missing glyph can be skipped; late
registrations refresh the relevant fallback chains.

This is different from the soft declaration fallback used when a font
itself cannot be resolved. For non-CJK scripts IMPE deliberately does not
replace missing glyphs one codepoint at a time, which could disrupt shaping
and word structure. Ordinary Unicode text uses fontspec/OpenType shaping
and the chosen script-specific behavior.

\subsection{Unicode-range and CJK routing}

A script-specific global binding can claim selected Unicode blocks while
unrelated text stays with its existing font. Existing local commands take
precedence within their argument and suppress the outer automatic range
router there. Leaving their scope restores global routing.

For CJK, IMPE uses xeCJK main/sans/mono channels. Shared Han characters
retain the active document CJK selection, while kana and Hangul can be
claimed as distinct script ranges. Use local selection such as
\verb|\JP{...}| or \verb|\KR{...}| when a regional form of a shared
Han ideograph must be chosen explicitly. Unified conflict resolution for
new, overlapping global Writing Profiles is not a language-detection
engine; the existing routing and arbitration limits still apply.

\subsection{Shaping, line breaking and writing directions}

OpenType shaping for Indic, Arabic and other complex scripts remains the
responsibility of \XeTeX{} and the selected OpenType fonts. IMPE supplies
script, language and feature settings and selected script-specific breaking
behavior: Sanskrit Devanagari akshara boundaries, Tibetan tsheg boundaries
and Thai dictionary-based line breaks are supported in their documented
routes. Right-to-left behavior is attached to appropriate presets and
renderers. The built-in local rotation route provides Mongolian, Manchu
and Old Uyghur-style \emph{local} vertical forms; it is not a general
four-direction page-layout engine.

\subsection{XSR 1.3: Egyptian, Khitan and inline glyph images}

XSR is vendored with IMPE rather than installed as a separate IMPE feature
named \texttt{xsr}. Two font modules use its active Unicode detector and
script backends; the independent \texttt{glyphs} feature uses its inline
vector/image backend.

\begin{description}
  \item[Egyptian hieroglyphs] Load \verb|\UseFont{egyptian}|, then use the
    renderer's supported Unicode hieroglyphic controls or
    \verb|\IMPEEgyptianText{...}| where explicit selection is preferable.
  \item[Khitan Small Script] Load \verb|\UseFont{khitan_small}| and use
    \verb|\KHS{...}|, the default \emph{native XSR stack} renderer.
    \verb|\KHSstack{...}| and \verb|\KHSstackblock{...}| remain
    explicit aliases. Unicode cluster separators and U+16FE4 fillers
    are retained; XSR 1.3 uses normal host-text interword spacing and
    vertically centers stacks of differing heights.
  \item[Registered inline images] Load \verb|\UseFeature{glyphs}| and
    register glyph assets with \verb|\IMPEGlyphRegister|; insert them
    through \verb|\IMPEGlyph|. This feature does not install the
    Egyptian or Khitan detector merely by being loaded.
\end{description}

For dynamic rendering, provide Python 3.11+, the installed private XSR
runtime or equivalent configured renderer, XeLaTeX shell escape (or XSR
preprocessing), and readable source/font files. Do not attempt
\verb|\UseFeature{xsr}|: there is no such feature identifier.

XSR 1.1 and later can preserve the original linear Unicode source of
Egyptian quadrats and Khitan clusters via PDF \texttt{/ActualText},
including controls and separators, while drawing the stacked visual output.
This is not an invisible duplicate text layer or a guarantee that every
PDF reader copies identically. IMPE does not claim general LaTeX-source
recovery for equations, tables or arbitrary image glyphs.

\subsection{Further font documentation}

The complete family catalog, source/font resolution, writing options,
style-face mapping, compatibility presets, XSR requirements and
script-specific notes are maintained in \texttt{docs/FONTS.md} and
\texttt{docs/USEWRITING-1.0.5.md}. The public Showcase illustrates
the actual rendered output.

\section{Layouts and Features}

\subsection{Layout presets}

Layouts describe document-level presentation.  A public layout preset is a
structured bundle of page geometry, text spacing, header behavior, book
behavior, or slide behavior.  The individual low-level components are internal;
normal documents select the public preset as a whole.

The public presets in IMPE 1.x are:

\begin{description}
  \item[\texttt{en\_doc}] English article/report layout.
  \item[\texttt{zh\_doc}] Chinese article/report layout.
  \item[\texttt{en\_book}] English book/report layout with book-oriented page
    and running-head behavior.
  \item[\texttt{zh\_book}] Chinese book/report counterpart.
  \item[\texttt{beamer}] presentation layout for Beamer documents.
\end{description}

Presets declare which underlying document classes they support.  IMPE checks
that compatibility before applying the preset and reports mismatched document
classes.

Wrapper classes select the corresponding layout automatically.  Explicit
selection is useful with raw standard classes or when a report should use a
book-oriented preset:

\begin{verbatim}
\UseLayout{en_book}
\end{verbatim}

The exact component values and geometry are documented in
\texttt{docs/LAYOUTS.md}.

\subsection{Feature modules}

Features are flat, composable modules. Each feature has its own
second preset layer.  A feature id simply loads the corresponding optional
functionality the first time it is requested; repeated requests are ignored.

The public feature set currently consists of:

\begin{description}
  \item[\texttt{math}] standard mathematics packages, theorem environments,
    and optional math-font selection.
  \item[\texttt{hyperlinks}] hyperlinks, PDF bookmarks, stable destinations,
    and repository hyperlink defaults.
  \item[\texttt{citations}] \texttt{biblatex}/\texttt{biber} citation support
    with selectable citation-style presets.
  \item[\texttt{index}] index generation and IMPE term/index helpers.
  \item[\texttt{tables}] table packages, reusable column types, and table
    environments.
  \item[\texttt{image}] figure, image, and multi-panel helpers.
  \item[\texttt{lists\_envs}] reusable example-block formatting.
  \item[\texttt{headers}] configurable running heads for article-, report-,
    and book-like documents.
\end{description}

A document loads any combination it requires:

\begin{verbatim}
\UseFeatures{math,hyperlinks,tables,image}
\end{verbatim}

or

\begin{verbatim}
\UseTemplateSet{
  features = {citations,index,hyperlinks}
}
\end{verbatim}

\subsection{Mathematics}

The \texttt{math} feature loads the standard AMS-oriented mathematics stack and
defines theorem-like environments including \texttt{theorem}, \texttt{lemma},
\texttt{proposition}, \texttt{corollary}, \texttt{definition},
\texttt{example}, and \texttt{remark}. Computer Modern mathematics remains the
default when a text family is loaded; an explicit math route selects another
math alphabet.

A different math family can be selected before the feature is loaded.  For
example,

\begin{verbatim}
\UseMathFont{libertinus}
\UseFeature{math}
\end{verbatim}

uses Libertinus Math through \texttt{unicode-math}.  Other supported routes and
the exact package stack are listed in \texttt{docs/FEATURES.md}.

\subsection{Citations, indexes, tables, and figures}

The feature modules provide useful defaults while retaining access to the
underlying \LaTeX{} ecosystem. The \texttt{citations} feature, for example, uses
\texttt{biblatex} and \texttt{biber}; documents still use ordinary commands
such as \verb|\addbibresource|, \verb|\textcite|, and
\verb|\printbibliography|.  Citation presets such as APA, GB/T 7714, numeric,
or author--year may be selected before the feature is loaded.

The \texttt{drawing} feature loads TikZ/PGF, pgfplots (compatibility 1.18), and
forest for their ordinary diagram, plot and tree syntax.

The \texttt{index} feature loads hyperlinks, builds on \texttt{imakeidx} and adds the
\verb|\Term| helper for typesetting and indexing a term together. Its printed
index page number links to the recorded occurrence even after page-number
resets. First-occurrence-per-key semantics are unchanged. For manual processing,
run \texttt{texindy -L english -C utf8 -M <job>-impe <job>.idx} in the output
directory, then rerun XeLaTeX twice. The
\texttt{tables} and \texttt{image} features similarly provide convenient
higher-level environments while leaving native \texttt{booktabs},
\texttt{graphicx}, and related commands available.

Throughout the feature system, IMPE reduces repetitive setup while preserving
ordinary \LaTeX{} commands.

\subsection{Headers and hyperlinks}

The \texttt{headers} feature supplies running heads through
\texttt{fancyhdr}.  By default it derives a fixed running title from the first
line of \verb|\title|; \verb|\HeaderTitle{...}| can be used when a shorter form
is preferable.  Article-like documents use section information for the moving
head, while report- and book-like documents use chapter information.

The \texttt{hyperlinks} feature loads \texttt{hyperref} and
\texttt{bookmark} with Unicode-aware PDF settings, hidden link borders, linked
table-of-contents structures, and stable destinations.  Documents may still use
ordinary \verb|\label|, \verb|\ref|, \verb|\href|, and \verb|\url| commands.

For the full list of feature-specific commands and environments, consult
\texttt{docs/FEATURES.md}; this section introduces their main uses.

\section{Compatibility and Migration}\label{sec:compatibility}

IMPE 1.x establishes the \texttt{impe*} names as the canonical public
interface.  New documents should therefore use \texttt{impe.sty},
\texttt{impeart}, \texttt{impebook}, \texttt{impereport},
\texttt{impebeamer}, and their Chinese variants.

Earlier development releases used the \texttt{next*} naming scheme.  Those
entry points remain supported compatibility wrappers in 1.x, with their
existing status maintained. Existing documents that load the
\texttt{nextart} class or \texttt{nextsystem} package continue to compile
against the same underlying implementation.
The compatibility layer preserves old documents; canonical \texttt{impe*}
entries serve new work.

The direct installer also recognizes a managed installation from the 0.1.x
series under \texttt{tex/latex/nextsystem/}.  During migration it installs the
canonical runtime under \texttt{tex/latex/impe/}, migrates recognized local
override files when appropriate, removes files known to belong to the old
managed runtime, and leaves unrelated user-created files or subdirectories in
place.

A legacy \texttt{nextsystem.local.tex} can therefore survive migration as a
local override. New local configuration should use \texttt{impe.local.tex}.

Documents written with old wrapper names continue to compile through the
compatibility runtime. Authors may update those names when convenient.

\section{Package Structure and Extension Points}

The repository separates public entry points, stable mechanisms, registrations,
specialized implementations, and local resources:

\begin{verbatim}
package/   public package and class entry points
core/      stable subsystem mechanisms
catalog/   public ids and registration data
modules/   extendable or specialized implementations
assets/    local runtime resources such as fonts
scripts/   installation and release tooling
manual/    manual sources
docs/      detailed subsystem documentation
\end{verbatim}

\subsection{The stable core}

The \texttt{core/} layer contains stable shared mechanisms for new font
families, layout presets, and features. It implements
the font registry and routing machinery, the layout registry and compatibility
checks, the feature loader, and related shared behavior.

This distinction matters for maintenance.  If adding one ordinary OpenType
font requires editing the font-routing engine, or adding one feature requires
changing the generic feature loader, the extension model has failed.  The core
is therefore intentionally separated from the growing catalog of concrete
resources.

\subsection{Catalogs}

The \texttt{catalog/} layer assigns public ids and metadata.  It contains the
central registrations for fonts, layouts, and features.  For many extensions,
this is the correct place to begin.

An ordinary font family using standard \texttt{fontspec} shaping can
usually be expressed entirely as a catalog registration.  Likewise, a public
layout preset is assembled from existing internal components, and a feature id
maps to its implementation module.  Keeping those declarations centralized
makes the public surface auditable in one place.

\subsection{Modules}

The \texttt{modules/} layer is reserved for functionality that genuinely needs
its own implementation: feature modules, specialized script support, or
reusable internal layout components.  A new module should therefore represent a
real behavioral extension; font names belong in catalog registrations.

This rule is particularly important for multilingual support.  Most scripts
should use ordinary OpenType shaping and catalog metadata.  Specialized modules
serve writing systems or rendering strategies requiring behavior beyond the
generic core.

\subsection{Extending IMPE}

Developers extending the system should begin with the detailed registration
interfaces in
\texttt{docs/SYSTEM.md} and then consult the relevant subsystem document:
\texttt{docs/FONTS.md}, \texttt{docs/LAYOUTS.md}, or
\texttt{docs/FEATURES.md}.

The general rule is simple: prefer a registration over a new mechanism, and
prefer a specialized module over a change to the stable core when the behavior
is genuinely exceptional.

\section{Further Documentation and Feedback}

This manual is the normal starting point for using IMPE.  The repository
contains more detailed technical documentation for users who need a complete
reference or who want to extend the system.

\begin{description}
  \item[\texttt{docs/SYSTEM.md}] architecture, public entry layers, template
    setup, release model, and extension boundaries.
  \item[\texttt{docs/FONTS.md}] family registration, local and global modes,
    routing, fallback, shaping, writing behavior, vertical support, and
    XSR integration.
  \item[\texttt{docs/LAYOUTS.md}] public layout presets, compatibility targets,
    and the component model from which presets are built.
  \item[\texttt{docs/FEATURES.md}] public feature ids and the commands,
    environments, package stacks, and compilation requirements supplied by each
    module.
  \item[\texttt{CHANGELOG.md}] changes in published versions.
  \item[Showcase] the full visual showcase is available from the
    \href{\IMPEShowcaseURL}{project repository}.
\end{description}

The project repository is available at

\begin{center}
\href{\IMPERepo}{\texttt{github.com/KelvinYangBK67/IMPE-LaTeX-System}}
\end{center}

Bug reports, feature requests, documentation corrections, portability reports,
and other concrete feedback are welcome through GitHub Issues:

\begin{center}
\href{\IMPEIssues}{\texttt{github.com/KelvinYangBK67/IMPE-LaTeX-System/issues}}
\end{center}

When reporting a typesetting problem, it is useful to include a minimal source
file, the \TeX{} engine and \TeX{} Live version, the IMPE version, and the relevant
font information.  Reports concerning a local or privately supplied font should
also distinguish between a font-coverage problem and an IMPE routing problem
where possible.

The English manual is the authoritative manual text.  Other language editions
may adapt prose and terminology to their own technical conventions, but their
commands, examples, and description of IMPE behavior should remain synchronized
with this edition.



\appendix

\section{Quick Reference}
\label{app:quick-reference}

This appendix summarizes the documented public interface of IMPE 1.x.
It serves as a lookup companion to the explanations in the main text and
detailed subsystem documentation under \texttt{docs/}. Internal implementation
commands, including those beginning with \texttt{\string\Next...}, are
documented with their implementation.

\subsection{Recommended font-loading model}

For ordinary use, prefer \texttt{fonts = \{...\}} inside
\texttt{\string\UseTemplateSet} or use \texttt{\string\UseFont\{...\}}
directly. Omitting the mode makes IMPE follow the registered behavior of the
family: it may define a local command, activate a document-wide font, or install
script- or range-specific automatic routing as appropriate.

Use an explicit mode to override that automatic behavior deliberately:

\begin{verbatim}
\UseFont{sanskrit}          % recommended: registered automatic behavior
\UseFont{sanskrit}[local]   % explicitly local only
\UseFont{sanskrit}[global]  % explicitly global/range-global

\UseTemplateSet{
  fonts = {sanskrit,tibetan,arabic},
  features = {math,hyperlinks}
}
\end{verbatim}

The convenience commands \texttt{\string\UseGlobalFont(s)} and
\texttt{\string\UseLocalFont(s)} express the same explicit distinction.
The \texttt{globalfonts} and \texttt{mainfonts} template keys are likewise
explicit global overrides; \texttt{fonts} provides automatic selection.

\subsection{Document classes and package entry points}

\begin{longtable}{@{}p{0.31\linewidth}p{0.63\linewidth}@{}}
\toprule
\textbf{Entry point} & \textbf{Purpose} \\
\midrule
\endfirsthead
\toprule
\textbf{Entry point} & \textbf{Purpose} \\
\midrule
\endhead
\texttt{impeart} &
English article wrapper; defaults to \texttt{en\_doc} with CMU as the base global font. \\
\texttt{impeart\_zh} &
Chinese article wrapper; defaults to \texttt{zh\_doc} with CMU and Shanggu. \\
\texttt{impereport} &
English report wrapper; defaults to \texttt{en\_doc}. \\
\texttt{impereport\_zh} &
Chinese report wrapper; defaults to \texttt{zh\_doc}. \\
\texttt{impebook} &
English book wrapper; defaults to \texttt{en\_book}. \\
\texttt{impebook\_zh} &
Chinese book wrapper; defaults to \texttt{zh\_book}. \\
\texttt{impebeamer} &
English Beamer wrapper; defaults to the \texttt{beamer} layout. \\
\texttt{impebeamer\_zh} &
Chinese Beamer wrapper; defaults to the \texttt{beamer} layout with CMU and Shanggu. \\
\texttt{impe.sty} &
Loads IMPE under an ordinary \LaTeX{} class chosen by the document. \\
\bottomrule
\end{longtable}

The legacy \texttt{next*} package and class names remain compatibility aliases
in version 1.x, but new documents should use the canonical \texttt{impe*}
entry points.

\subsection{Core setup commands}

\begin{longtable}{@{}p{0.38\linewidth}p{0.56\linewidth}@{}}
\toprule
\textbf{Command} & \textbf{Purpose} \\
\midrule
\endfirsthead
\toprule
\textbf{Command} & \textbf{Purpose} \\
\midrule
\endhead
\texttt{\string\UseTemplateSet\{...\}} &
Configure layout, fonts, and features in one declaration. \\
\texttt{\string\UseFont\{id\}[mode]} &
Load one registered font family; omit \texttt{mode} for the recommended automatic behavior, or use \texttt{local}/\texttt{global} explicitly. \\
\texttt{\string\UseFonts\{ids\}[mode]} &
Load several registered font families with the same optional mode. \\
\texttt{\string\UseGlobalFont\{id\}} &
Explicitly request the global mode of one family. \\
\texttt{\string\UseGlobalFonts\{ids\}} &
Explicitly request the global mode of several families. \\
\texttt{\string\UseMainFont\{id\}} &
Alias of \texttt{\string\UseGlobalFont}; retained as a document-oriented name. \\
\texttt{\string\UseMainFonts\{ids\}} &
Alias of \texttt{\string\UseGlobalFonts}. \\
\texttt{\string\UseLocalFont\{id\}} &
Explicitly request the local mode of one family. \\
\texttt{\string\UseLocalFonts\{ids\}} &
Explicitly request the local mode of several families. \\
\texttt{\string\UseLayout\{id\}} &
Load one registered layout preset. \\
\texttt{\string\UseLayouts\{ids\}} &
Load several registered layout presets. \\
\texttt{\string\UseFeature\{id\}} &
Load one optional feature module. \\
\texttt{\string\UseFeatures\{ids\}} &
Load several optional feature modules. \\
\texttt{\string\SetCatalogFontRoot\{path\}} &
Override the root directory from which catalogued font files are resolved. \\
\texttt{\string\subtitle\{text\}} &
Set the document subtitle used by the IMPE title block. \\
\texttt{\string\HeaderTitle\{text\}} &
Override the fixed title used by the \texttt{headers} feature. \\
\texttt{\string\HeaderStyle\{style\}} &
Choose the header style; documented values are \texttt{running} and \texttt{title}. \\
\bottomrule
\end{longtable}

\subsection{Template-set keys}

\begin{longtable}{@{}p{0.23\linewidth}p{0.69\linewidth}@{}}
\toprule
\textbf{Key} & \textbf{Meaning} \\
\midrule
\endfirsthead
\toprule
\textbf{Key} & \textbf{Meaning} \\
\midrule
\endhead
\texttt{layout} &
Select one public layout preset. \\
\texttt{fonts} &
Load one or more registered families using their automatic registered behavior; this is the recommended key for normal font loading. \\
\texttt{globalfonts} &
Explicitly force the listed families through their global mode or global routing path. \\
\texttt{mainfonts} &
Alias of \texttt{globalfonts}. \\
\texttt{features} &
Load one or more optional feature modules. \\
\bottomrule
\end{longtable}

\subsection{Public layout identifiers}

\begin{longtable}{@{}p{0.23\linewidth}p{0.69\linewidth}@{}}
\toprule
\textbf{ID} & \textbf{Purpose} \\
\midrule
\endfirsthead
\toprule
\textbf{ID} & \textbf{Purpose} \\
\midrule
\endhead
\texttt{en\_doc} &
English article/report layout with A4 geometry and English text spacing. \\
\texttt{zh\_doc} &
Chinese article/report layout with A4 geometry and Chinese text spacing. \\
\texttt{report} &
Report-only compatibility preset with the same page and text components as \texttt{zh\_doc}. \\
\texttt{en\_book} &
English book/report layout with book geometry, running heads, and right-hand chapter openings. \\
\texttt{zh\_book} &
Chinese book/report layout with book geometry, Chinese text spacing, running heads, and right-hand chapter openings. \\
\texttt{beamer} &
Presentation layout with the IMPE Beamer typography and navigation setup. \\
\bottomrule
\end{longtable}

\subsection{Public feature identifiers}

\begin{longtable}{@{}p{0.23\linewidth}p{0.69\linewidth}@{}}
\toprule
\textbf{ID} & \textbf{Purpose} \\
\midrule
\endfirsthead
\toprule
\textbf{ID} & \textbf{Purpose} \\
\midrule
\endhead
\texttt{math} &
Load the IMPE mathematics stack and theorem environments. \\
\texttt{hyperlinks} &
Load hyperlinking, PDF bookmarks, linked headings, and bidirectional footnote links. \\
\texttt{citations} &
Load \texttt{biblatex}/\texttt{csquotes} citation support and IMPE citation-style helpers. \\
\texttt{index} &
Load index support and the \texttt{\string\Term} helper. \\
\texttt{tables} &
Load the table stack, column types, and IMPE table environments. \\
\texttt{image} &
Load image, caption, and panel-figure helpers. \\
\texttt{lists\_envs} &
Load the \texttt{ExampleBlock} display environment. \\
\texttt{headers} &
Load configurable running heads for article-, report-, and book-like documents. \\
\texttt{bib}, \texttt{header} &
Compatibility aliases for \texttt{citations} and \texttt{headers}. \\
\bottomrule
\end{longtable}

\subsection{Feature-specific commands and values}

\begin{longtable}{@{}p{0.51\linewidth}p{0.43\linewidth}@{}}
\toprule
\textbf{Command or setting} & \textbf{Purpose / documented values} \\
\midrule
\endfirsthead
\toprule
\textbf{Command or setting} & \textbf{Purpose / documented values} \\
\midrule
\endhead
\texttt{\string\UseMathFont\{value\}} &
Choose the mathematics font route before loading \texttt{math}: \texttt{auto}, \texttt{libertinus}, \texttt{newcm}, \texttt{mlmodern}, or another font name passed to \texttt{\string\setmathfont}. \\
\texttt{\string\UseCitationStyle\{value\}} &
Choose a citation preset before loading \texttt{citations}: \texttt{APA}, \texttt{GB}, \texttt{numeric}, or \texttt{author-year}. \\
\texttt{\string\SetCitationBiblatexOptions\{...\}} &
Replace the effective \texttt{biblatex} option list before loading \texttt{citations}. \\
\texttt{\string\HeaderTitle\{text\}} &
Set a short fixed title for running heads. \\
\texttt{\string\HeaderStyle\{value\}} &
Select \texttt{running} or \texttt{title}. \\
\texttt{\string\IndexTitle} &
Optional index-title override, defined before loading \texttt{index}. \\
\texttt{\string\Term[options]\{display\}[description]} &
Print a bold term and index its first occurrence; options include \texttt{sort}, \texttt{key}, and \texttt{parentheses}. \\
\texttt{\string\TablesSetup} &
Apply or re-apply the IMPE table setup. \\
\texttt{\string\Panel} &
Insert one panel inside \texttt{PanelFigure} or \texttt{PanelFigure*}. \\
\texttt{\string\TemplateFigurePaths} &
Configure the default search paths used by the image feature. \\
\texttt{\string\OneImageDefaultWidth} &
Set the default width for \texttt{OneImage}. \\
\texttt{\string\OneImageMaxHeight} &
Set the maximum default image height for \texttt{OneImage}. \\
\texttt{\string\OneImageDefaultPlacement} &
Set the default float placement for \texttt{OneImage}. \\
\texttt{\string\PanelDefaultCols} &
Set the default number of columns used by panel figures. \\
\texttt{\string\PanelDefaultHeight} &
Set the default panel height. \\
\texttt{\string\PanelDefaultMode} &
Set the default panel layout mode. \\
\texttt{\string\PanelDefaultPlacement} &
Set the default float placement for panel figures. \\
\bottomrule
\end{longtable}

Standard commands provided directly by packages loaded by a feature---for
example \texttt{\string\addbibresource}, \texttt{\string\printbibliography},
\texttt{\string\printindex}, or the \texttt{booktabs} rule commands---retain
their normal package syntax as provided by \LaTeX{}.

\subsection{Public environments and table column types}

\begin{longtable}{@{}p{0.37\linewidth}p{0.57\linewidth}@{}}
\toprule
\textbf{Name} & \textbf{Purpose} \\
\midrule
\endfirsthead
\toprule
\textbf{Name} & \textbf{Purpose} \\
\midrule
\endhead
\texttt{theorem}, \texttt{lemma}, \texttt{proposition}, \texttt{corollary} &
Numbered theorem-like environments provided by the \texttt{math} feature. \\
\texttt{definition}, \texttt{example}, \texttt{remark} &
Additional theorem-style environments provided by the \texttt{math} feature. \\
\texttt{TableInlineFit} &
Compact table environment intended to fit an inline table to the available width. \\
\texttt{TableLong} &
Long-table environment for material that may span pages. \\
\texttt{TableBook}, \texttt{TableBookX} &
Book-style table environments. \\
\texttt{TableBookNotes} &
Book-style table environment with notes. \\
\texttt{NiceBooktable}, \texttt{NiceBooktableX}, \texttt{NiceBooktableNotes} &
Alternative book-style table helpers provided by the \texttt{tables} feature. \\
\texttt{L}, \texttt{C}, \texttt{R} &
Ragged-right, centered, and ragged-left \texttt{tabularx} columns. \\
\texttt{P\{w\}}, \texttt{M\{w\}}, \texttt{B\{w\}} &
Fixed-width paragraph columns with ragged-right, centered, and ragged-left alignment. \\
\texttt{OneImage} &
Single-image figure; appears inline under Beamer. \\
\texttt{OneImageInline} &
Centered non-floating image. \\
\texttt{PanelFigure} &
Captioned multi-panel figure with subcaptions outside Beamer. \\
\texttt{PanelFigure*} &
Uncaptioned multi-panel figure. \\
\texttt{ExampleBlock} &
Indented italic display block for examples, quotations, linguistic data, or teaching material. \\
\bottomrule
\end{longtable}

\subsection{Registered font-family identifiers}

The table below lists the font-family IDs accepted by
\texttt{\string\UseFont} in IMPE 1.x.  The \textbf{Auto} column summarizes
what a mode-less \texttt{\string\UseFont\{id\}} normally activates:

\begin{description}
  \item[\texttt{G}] document-global family;
  \item[\texttt{L}] local family/command;
  \item[\texttt{L+G}] local access plus full document-global activation;
  \item[\texttt{L+R}] local access plus automatic script- or range-limited global routing.
\end{description}

The automatic route is the recommended interface.  An explicit
\texttt{[local]} or \texttt{[global]} request deliberately bypasses part of
that automatic decision and should therefore be used only when required.

\begin{longtable}{@{}p{0.30\linewidth}p{0.11\linewidth}p{0.18\linewidth}p{0.31\linewidth}@{}}
\toprule
\textbf{Family ID} & \textbf{Auto} & \textbf{Local command} & \textbf{Typical role} \\
\midrule
\endfirsthead
\toprule
\textbf{Family ID} & \textbf{Auto} & \textbf{Local command} & \textbf{Typical role} \\
\midrule
\endhead
\texttt{cmu} & G & --- & Computer Modern Unicode base family. \\
\texttt{noto} & L+G & \texttt{\string\NOT} & General Noto Latin family. \\
\texttt{times} & L+G & \texttt{\string\TIM} & Times/Arial/Consolas system-family route. \\
\texttt{gentium} & L+G & \texttt{\string\GEN} & Gentium Plus. \\
\texttt{charis} & L+G & \texttt{\string\CHA} & Charis SIL. \\
\texttt{libertinus} & L+G & \texttt{\string\LIB} & Libertinus serif/sans/mono family. \\
\texttt{mlmodern} & L & \texttt{\string\MLM} & Latin Modern / MLModern local route. \\

\texttt{anatolian} & L & \texttt{\string\CA} & Carian / Anatolian script material. \\
\texttt{coptic} & L & \texttt{\string\CO} & Coptic. \\
\texttt{bopomofo} & L & \texttt{\string\ZY} & Bopomofo. \\
\texttt{cuneiform} & L & \texttt{\string\CU} & Cuneiform. \\
\texttt{glagolitic} & L & \texttt{\string\GL} & Glagolitic. \\
\texttt{italic} & L & \texttt{\string\OI} & Old Italic. \\
\texttt{hungarian} & L & \texttt{\string\OH} & Old Hungarian. \\
\texttt{runic} & L & \texttt{\string\RU} & Runic. \\
\texttt{armenian} & L & \texttt{\string\HY} & Armenian. \\
\texttt{georgian} & L & \texttt{\string\KA} & Georgian. \\

\texttt{hindi} & L+R & \texttt{\string\HI} & Hindi Devanagari with registered range routing. \\
\texttt{sanskrit} & L+R & \texttt{\string\SA} & Sanskrit Devanagari with Sanskrit-aware routing. \\
\texttt{devanagari} & L & \texttt{\string\DEV} & Generic Devanagari shaping. \\
\texttt{tamil} & L & \texttt{\string\TA} & Tamil. \\
\texttt{brahmi} & L & \texttt{\string\BR} & Brahmi. \\
\texttt{tibetan} & L+R & \texttt{\string\TI} & Tibetan with tsheg-aware line-breaking behavior. \\
\texttt{segoe} & L & \texttt{\string\SEG} & Segoe Historic local route. \\
\texttt{thai} & L & \texttt{\string\TH} & Thai. \\

\texttt{arabic} & L & \texttt{\string\AR} & Arabic; serif/sans and alternate stylistic faces. \\
\texttt{urdu} & L & \texttt{\string\UR} & Urdu Nastaliq. \\
\texttt{aramaic} & L & \texttt{\string\IA} & Imperial Aramaic. \\
\texttt{nabataean} & L & \texttt{\string\NB} & Nabataean. \\
\texttt{hebrew} & L & \texttt{\string\HE} & Hebrew. \\
\texttt{syriac} & L & \texttt{\string\SY} & Syriac. \\
\texttt{syriac\_eastern} & L+R & \texttt{\string\SYE} & Eastern Syriac with registered range routing. \\
\texttt{kharosthi} & L & \texttt{\string\KH} & Kharosthi. \\
\texttt{khitan\_small} & L & \texttt{\string\KHS} & Khitan Small Script with its specialized composer. \\
\texttt{pahlavi\_parthian} & L+R & \texttt{\string\PAR} & Inscriptional Parthian. \\
\texttt{pahlavi\_inscriptional} & L+R & \texttt{\string\PAH} & Inscriptional Pahlavi. \\
\texttt{pahlavi\_psalter} & L & \texttt{\string\PSP} & Psalter Pahlavi with specialized support. \\
\texttt{avestan} & L & \texttt{\string\AV} & Avestan. \\
\texttt{manichaean} & L & \texttt{\string\MA} & Manichaean. \\
\texttt{phoenician} & L & \texttt{\string\PH} & Phoenician. \\
\texttt{samaritan} & L & \texttt{\string\SM} & Samaritan. \\
\texttt{sogdian} & L & \texttt{\string\SG} & Sogdian. \\
\texttt{sogdian\_old} & L+R & \texttt{\string\SGO} & Old Sogdian with registered range routing. \\

\texttt{chinese\_simplified} & L+R & \texttt{\string\SC} & Simplified Chinese CJK family. \\
\texttt{chinese\_traditional} & L+R & \texttt{\string\TC} & Traditional Chinese CJK family. \\
\texttt{japanese} & L+R & \texttt{\string\JP} & Japanese; automatic routing preserves shared Han priority. \\
\texttt{wenjin} & L+R & \texttt{\string\WJ} & WenJin Mincho with CJK fallback chain. \\
\texttt{shanggu} & G & --- & Traditional-Chinese-oriented global CJK family. \\
\texttt{sim} & G & --- & Simplified-Chinese-oriented system CJK family. \\
\texttt{korean} & L+R & \texttt{\string\KR} & Korean/Hangul family. \\
\texttt{tangut} & L+R & \texttt{\string\TG} & Tangut with CJK-aware automatic routing. \\

\texttt{mongolian} & L & \texttt{\string\MO} & Mongolian with built-in vertical layout route. \\
\texttt{mongolian\_baiti} & L & \texttt{\string\MOb} & Mongolian Baiti with vertical layout route. \\
\texttt{manchu} & L & \texttt{\string\MC} & Manchu with vertical layout route. \\
\texttt{turkic} & L & \texttt{\string\OT} & Old Turkic. \\
\texttt{uyghur} & L & \texttt{\string\UY} & Old Uyghur with vertical/rotated layout route. \\

\texttt{vietnamese\_quocngu} & L & \texttt{\string\VI} & Vietnamese Quốc Ngữ. \\
\texttt{vietnamese\_hannom} & L & \texttt{\string\HN} & Vietnamese Hán-Nôm / CJK text. \\
\bottomrule
\end{longtable}

\subsection{Advanced extension interfaces}

These commands support extension of IMPE.

\begin{longtable}{@{}p{0.39\linewidth}p{0.55\linewidth}@{}}
\toprule
\textbf{Command} & \textbf{Purpose} \\
\midrule
\endfirsthead
\toprule
\textbf{Command} & \textbf{Purpose} \\
\midrule
\endhead
\texttt{\string\FontRegisterFamily\{...\}} &
Register a new font family and its local/global metadata. \\
\texttt{\string\LayoutPresetRegister\{...\}} &
Register a structured layout preset from internal layout components. \\
\texttt{\string\LayoutPresetDeclare\{...\}} &
Lower-level layout-preset declaration interface used when building the system. \\
\bottomrule
\end{longtable}

For the full registration syntax and extension rules, see
\texttt{docs/FONTS.md}, \texttt{docs/LAYOUTS.md}, and
\texttt{docs/SYSTEM.md}.

\section{Showcase}\label{app:showcase}

The following pages reproduce the complete canonical IMPE showcase.  The
showcase is a visual companion to the reference manual,
demonstrating multilingual font routing, complex shaping, right-to-left
scripts, vertical writing, CJK regional forms, and combinations of document
features.

The latest showcase can be obtained from the
\href{\IMPEShowcaseURL}{current main branch}; the edition belonging to this
manual is preserved by the
\href{\IMPETaggedShowcaseURL}{v\IMPEVersion{} tag}.

\IfFileExists{\IMPEShowcaseFile}
  {\includepdf[pages=1-3,pagecommand={}]{\IMPEShowcaseFile}%
   \includepdf[pages=4-6,landscape=true,pagecommand={}]{\IMPEShowcaseFile}%
   \includepdf[pages=7-,pagecommand={}]{\IMPEShowcaseFile}}
  {\PackageError{impe-manual}
     {Required repository resource ../showcase/impe-showcase.pdf is missing}
     {Compile from manual/en or run scripts/build_manual.ps1 from the repository root.}}

\end{document}
